\documentclass[UTF8,a4paper,11pt]{ctexart}

\usepackage{amsmath,amssymb,amsthm}
\usepackage{geometry}
\usepackage{hyperref}
\usepackage{enumitem}
\usepackage{booktabs}

\geometry{margin=1in}
\hypersetup{
  colorlinks=true,
  linkcolor=blue,
  citecolor=blue,
  urlcolor=blue
}

\newtheorem{definition}{定义}
\newtheorem{assumption}{假设}
\newtheorem{lemma}{引理}
\newtheorem{theorem}{定理}
\newtheorem{corollary}{推论}
\newtheorem{remark}{说明}
\renewcommand{\proofname}{证明}

\newcommand{\R}{\mathbb{R}}
\newcommand{\E}{\mathbb{E}}
\newcommand{\V}{\mathcal{V}}
\newcommand{\B}{\mathcal{B}}
\newcommand{\F}{\mathcal{F}}
\newcommand{\Gap}{\mathcal{G}}
\newcommand{\norm}[1]{\left\|#1\right\|}

\title{分位数回归盒对偶联邦原始--对偶混合梯度法的严格收敛阶数证明}
\author{项目组}
\date{\today}

\begin{document}

\maketitle

\begin{abstract}
本文给出分位数回归盒对偶联邦原始--对偶混合梯度法在一个明确随机联邦模型下的收敛阶数证明。
证明中显式定义客户端抽样、用户小批量抽样、缓存方向以及陈旧偏差，并给出依赖客户端参与数、小批量大小、客户端不均衡因子和最大缓存延迟的收敛界。
主要结论如下：在全客户端、全样本的确定性情形下，遍历原始--对偶间隙具有 $\mathcal{O}(1/T)$ 收敛阶数；在随机联邦、部分客户端参与和小批量更新下，即使不使用方差缩减，也可以得到 $\mathcal{O}(1/\sqrt{T})$ 的期望遍历间隙界。方差缩减可以降低方差常数，但不是得到该阶数的必要条件。
\end{abstract}

\section{问题、联邦划分与鞍点形式}

\begin{definition}[带惩罚的分位数回归问题]
给定数据 $\{(x_i,y_i)\}_{i=1}^n$，其中 $x_i\in\R^p$，考虑
\begin{equation}
    \min_{\beta\in\B}
    F(\beta)
    :=
    \frac{1}{n}\sum_{i=1}^n
    \rho_\tau(y_i-x_i^\top\beta)
    +
    r(\beta),
    \label{eq:primal}
\end{equation}
其中 $\B\subset\R^p$ 是闭凸有界集合，$r$ 是真闭凸惩罚函数，
\[
    \rho_\tau(u)=u\bigl(\tau-\mathbf{1}\{u<0\}\bigr),
    \qquad \tau\in(0,1),
\]
为分位数回归的检查损失。
\end{definition}

\begin{remark}
这里显式加入有界集合 $\B$，是为了使随机误差项中的内积具有有限上界。在实际算法中，这可以由投影、参数截断、强制惩罚项，或者有界水平集条件保证。若使用最小最大凹惩罚或平滑截断绝对偏差惩罚等非凸惩罚，则严格全局鞍点收敛结论需改写为驻点型结论；本文证明针对凸惩罚情形。
\end{remark}

数据被划分到 $m$ 个客户端：
\[
    \mathcal{D}_1,\ldots,\mathcal{D}_m,\qquad
    n_j=|\mathcal{D}_j|,\qquad
    q_j=\frac{n_j}{n},\qquad
    \sum_{j=1}^m q_j=1.
\]
客户端 $j$ 的设计矩阵和对偶变量分别记为 $X_j$ 与 $v_j$。

\begin{definition}[盒对偶表示]
检查损失满足芬歇尔盒对偶表示
\begin{equation}
    \rho_\tau(u)
    =
    \max_{v\in[\tau-1,\tau]} vu.
    \label{eq:box-dual}
\end{equation}
因此，问题 \eqref{eq:primal} 等价于凸凹鞍点问题
\begin{equation}
    \min_{\beta\in\B}
    \max_{v\in\V}
    \mathcal{L}(\beta,v)
    :=
    r(\beta)
    +
    \frac{1}{n}v^\top(y-X\beta),
    \qquad
    \V=[\tau-1,\tau]^n.
    \label{eq:saddle}
\end{equation}
\end{definition}

\begin{definition}[客户端局部方向与全量方向]
定义客户端 $j$ 的局部方向为
\begin{equation}
    d_j(v_j)
    :=
    \frac{1}{n_j}X_j^\top v_j.
    \label{eq:local-direction}
\end{equation}
全量原始变量更新方向为
\begin{equation}
    g(v)
    :=
    \frac{1}{n}X^\top v
    =
    \sum_{j=1}^m q_j d_j(v_j).
    \label{eq:full-direction}
\end{equation}
分位数回归盒对偶联邦原始--对偶混合梯度法在服务器端使用 $g(v)$ 或其随机联邦估计量更新 $\beta$。
\end{definition}

\section{基本假设与有界性引理}

\begin{assumption}[有界性与凸性]
\label{assump:bounded}
存在常数 $R,D_\beta,Y_\B<\infty$，使得：
\begin{enumerate}[leftmargin=2em]
    \item 对所有样本，$\norm{x_i}\le R$；
    \item 集合 $\B$ 的直径不超过 $D_\beta$，即任意 $\beta,\beta'\in\B$ 满足
    \[
        \norm{\beta-\beta'}\le D_\beta;
    \]
    \item 对所有 $\beta\in\B$，残差有界：
    \[
        |y_i-x_i^\top\beta|\le Y_\B;
    \]
    \item $r$ 是真闭凸函数，并且其近端映射在 $\B$ 上良定义。
\end{enumerate}
\end{assumption}

\begin{assumption}[鞍点存在]
\label{assump:saddle}
鞍点问题 \eqref{eq:saddle} 至少存在一个鞍点，记为
$(\beta^\star,v^\star)$。也就是说，对任意 $(\beta,v)\in\B\times\V$，
\[
    \mathcal{L}(\beta^\star,v)
    \le
    \mathcal{L}(\beta^\star,v^\star)
    \le
    \mathcal{L}(\beta,v^\star).
\]
在本文假设下，若 $\B$ 非空紧凸、$\V$ 非空紧凸，且 $r$ 在 $\B$ 上下半连续，
则该条件由极小极大定理保证。
\end{assumption}

\begin{lemma}[局部方向有界]
\label{lem:direction-bound}
令
\[
    V_\tau:=\max\{\tau,1-\tau\}.
\]
由于 $v_i\in[\tau-1,\tau]$，对任意客户端 $j$ 和任意 $v_j$，有
\begin{equation}
    \norm{d_j(v_j)}
    \le
    R V_\tau
    =:G.
    \label{eq:G-bound}
\end{equation}
\end{lemma}

\begin{proof}
由三角不等式和假设 \ref{assump:bounded}，
\[
    \norm{d_j(v_j)}
    =
    \left\|
    \frac{1}{n_j}\sum_{i\in\mathcal{D}_j}x_i v_i
    \right\|
    \le
    \frac{1}{n_j}
    \sum_{i\in\mathcal{D}_j}
    \norm{x_i}|v_i|
    \le
    RV_\tau.
\]
\end{proof}

\section{确定性全量方法的 \texorpdfstring{$\mathcal{O}(1/T)$}{一除以T} 阶}

先考虑没有任何随机性的基准情形：每一轮所有客户端都参与，并且每个客户端都刷新全部本地对偶变量。令
\[
    K_0=-X/n.
\]

\begin{assumption}[确定性步长条件]
\label{assump:det}
步长 $\eta,\sigma>0$ 满足
\begin{equation}
    \eta\sigma \norm{K_0}_2^2
    =
    \eta\sigma \frac{\norm{X}_2^2}{n^2}
    <1.
    \label{eq:det-step}
\end{equation}
\end{assumption}

\begin{theorem}[确定性全量方法的遍历收敛阶数]
\label{thm:det}
在假设 \ref{assump:bounded}、\ref{assump:saddle} 和 \ref{assump:det} 下，
令 $(\beta^t,v^t)$ 为全量确定性盒对偶原始--对偶混合梯度迭代。定义
\[
    \bar{\beta}^T=\frac{1}{T}\sum_{t=1}^T\beta^t,
    \qquad
    \bar{v}^T=\frac{1}{T}\sum_{t=1}^T v^t.
\]
若 $(\beta^\star,v^\star)$ 是 \eqref{eq:saddle} 的鞍点，则
\begin{equation}
    \Gap(\bar{\beta}^T,\bar{v}^T)
    :=
    \mathcal{L}(\bar{\beta}^T,v^\star)
    -
    \mathcal{L}(\beta^\star,\bar{v}^T)
    \le
    \frac{1}{T}
    \left(
        \frac{\norm{\beta^\star-\beta^0}^2}{2\eta}
        +
        \frac{\norm{v^\star-v^0}^2}{2\sigma}
    \right).
    \label{eq:det-rate}
\end{equation}
因此确定性全量方法的遍历原始--对偶间隙为 $\mathcal{O}(1/T)$。
\end{theorem}

\begin{proof}
问题 \eqref{eq:saddle} 是标准凸凹鞍点问题：
\[
    r(\beta)+\langle K_0\beta,v\rangle-h^\ast(v),
    \qquad
    h^\ast(v):=-v^\top y/n+\iota_{\V}(v),
\]
其中 $\iota_\V$ 是盒集合 $\V$ 的示性函数。对偶变量更新是 $h^\ast$ 的近端步骤，即投影到 $\V$；原始变量更新是 $r$ 的近端步骤。
由原始--对偶混合梯度法的一步不等式，在 \eqref{eq:det-step} 下，对任意 $(\beta,v)\in\B\times\V$ 有
\begin{align}
    \mathcal{L}(\beta^{t+1},v)
    -
    \mathcal{L}(\beta,v^{t+1})
    &\le
    \frac{1}{2\eta}
    \left(
        \norm{\beta-\beta^t}^2
        -
        \norm{\beta-\beta^{t+1}}^2
    \right)
    \nonumber\\
    &\quad+
    \frac{1}{2\sigma}
    \left(
        \norm{v-v^t}^2
        -
        \norm{v-v^{t+1}}^2
    \right).
    \label{eq:pdhg-one-step}
\end{align}
令 $(\beta,v)=(\beta^\star,v^\star)$，并对 $t=0,\ldots,T-1$ 求和，右端平方距离项望远镜相消。再利用 $\mathcal{L}$ 对 $\beta$ 的凸性和对 $v$ 的凹性，将平均间隙转化为遍历平均点处的间隙，即得到 \eqref{eq:det-rate}。
\end{proof}

\section{客户端抽样与小批量方差}

下面给出随机联邦方向估计的显式方差界。

\subsection{无缓存偏差时的随机方向估计}

每轮从分布 $\{\pi_j\}_{j=1}^m$ 中有放回抽取 $K$ 个客户端，记为
$J_{t,1},\ldots,J_{t,K}$，其中 $\pi_j>0$ 且 $\sum_j\pi_j=1$。
在客户端 $j$ 上，用大小为 $B_j$ 的用户小批量构造局部方向估计 $\widehat d_j^t$，满足
\begin{equation}
    \E[\widehat d_j^t\mid\F_t]=d_j(v_j^t),
    \qquad
    \E\left[
    \norm{\widehat d_j^t-d_j(v_j^t)}^2
    \mid\F_t
    \right]
    \le
    \frac{G^2}{B_j}.
    \label{eq:local-batch-var}
\end{equation}
其中 $G=RV_\tau$ 来自引理 \ref{lem:direction-bound}。

\begin{definition}[逆概率加权方向估计]
定义全局方向估计量
\begin{equation}
    \widehat g^t
    =
    \frac{1}{K}
    \sum_{\ell=1}^K
    \frac{q_{J_{t,\ell}}}{\pi_{J_{t,\ell}}}
    \widehat d_{J_{t,\ell}}^t.
    \label{eq:ht-estimator}
\end{equation}
该形式是标准逆概率加权估计，用于修正客户端非均匀抽样带来的偏差。
\end{definition}

\begin{lemma}[客户端抽样与小批量方差界]
\label{lem:variance}
在上述抽样模型下，
\begin{equation}
    \E[\widehat g^t\mid\F_t]
    =
    g(v^t),
    \label{eq:unbiased}
\end{equation}
且存在一个绝对常数 $c_0>0$，使得
\begin{equation}
    \E\left[
    \norm{\widehat g^t-g(v^t)}^2
    \mid\F_t
    \right]
    \le
    \frac{c_0G^2}{K}
    \sum_{j=1}^m
    \frac{q_j^2}{\pi_j}
    \left(
        1+\frac{1}{B_j}
    \right)
    =:\Sigma_{K,B}^2.
    \label{eq:variance-bound}
\end{equation}
\end{lemma}

\begin{proof}
无偏性由
\[
    \E\left[
    \frac{q_{J}}{\pi_J}\widehat d_J^t
    \mid\F_t
    \right]
    =
    \sum_{j=1}^m
    \pi_j
    \frac{q_j}{\pi_j}
    d_j(v_j^t)
    =
    \sum_{j=1}^m q_j d_j(v_j^t)
    =
    g(v^t)
\]
得到。由于 $K$ 次抽样条件独立，
\[
    \E\left[
    \norm{\widehat g^t-g(v^t)}^2
    \mid\F_t
    \right]
    \le
    \frac{1}{K}
    \E\left[
    \left.
    \left\|
    \frac{q_J}{\pi_J}\widehat d_J^t
    \right\|^2
    \right|\F_t
    \right].
\]
又有
\[
    \E[\norm{\widehat d_j^t}^2\mid\F_t]
    \le
    2\norm{d_j(v_j^t)}^2
    +
    2\E[\norm{\widehat d_j^t-d_j(v_j^t)}^2\mid\F_t]
    \le
    2G^2\left(1+\frac{1}{B_j}\right).
\]
代入上式即可得到 \eqref{eq:variance-bound}，其中绝对常数可取 $c_0=2$。
\end{proof}

\begin{remark}[客户端不均衡因子]
若 $\pi_j=1/m$，则
\[
    \sum_{j=1}^m\frac{q_j^2}{\pi_j}
    =
    m\sum_{j=1}^m q_j^2.
\]
当客户端样本量均衡时，$q_j=1/m$，该项等于 $1$；当客户端样本量高度不均衡时，该项变大。因此，客户端样本量不均衡会放大随机方向方差。
\end{remark}

\section{缓存方向的陈旧偏差}

实际的缓存聚合并不总是使用当前 $d_j(v_j^t)$，而是使用客户端上次被刷新时的方向。
记客户端 $j$ 在第 $t$ 轮服务器可见的缓存时间为 $\tau_j(t)\le t$，缓存方向为
\[
    \widetilde d_j^t
    :=
    d_j(v_j^{\tau_j(t)}).
\]
缓存全量聚合方向为
\[
    \widetilde g^t
    =
    \sum_{j=1}^m q_j \widetilde d_j^t.
\]

\begin{assumption}[有界缓存延迟]
\label{assump:delay}
存在整数 $A\ge 0$，使得
\[
    0\le t-\tau_j(t)\le A,
    \qquad
    \text{对所有 }j,t\text{ 成立}.
\]
此外，对偶变量步长 $\sigma_t$ 非增。
\end{assumption}

\begin{lemma}[陈旧偏差显式界]
\label{lem:stale}
在假设 \ref{assump:bounded} 和 \ref{assump:delay} 下，
\begin{equation}
    \norm{\widetilde g^t-g(v^t)}
    \le
    RY_\B
    \sum_{s=t-A}^{t-1}\sigma_s
    \le
    RY_\B A\sigma_{t-A},
    \label{eq:stale-bound}
\end{equation}
其中当 $t-A<0$ 时求和从 $0$ 开始。
\end{lemma}

\begin{proof}
对偶变量的盒投影更新满足
\[
    v_i^{s+1}
    =
    \Pi_{[\tau-1,\tau]}
    \left(
        v_i^s+\sigma_s(y_i-x_i^\top\bar\beta^s)
    \right),
\]
或其小批量分块版本。投影算子是非扩张的，因此对任意被更新的坐标，
\[
    |v_i^{s+1}-v_i^s|
    \le
    \sigma_s |y_i-x_i^\top\bar\beta^s|
    \le
    \sigma_s Y_\B.
\]
未被更新的坐标变化为 $0$，所以上式仍成立。于是对客户端 $j$，
\begin{align*}
    \norm{d_j(v_j^t)-d_j(v_j^{\tau_j(t)})}
    &=
    \left\|
    \frac{1}{n_j}
    X_j^\top
    (v_j^t-v_j^{\tau_j(t)})
    \right\|  \\
    &\le
    \frac{1}{n_j}
    \sum_{i\in\mathcal{D}_j}
    \norm{x_i}
    |v_i^t-v_i^{\tau_j(t)}| \\
    &\le
    R\sum_{s=\tau_j(t)}^{t-1}\sigma_s Y_\B
    \le
    RY_\B\sum_{s=t-A}^{t-1}\sigma_s .
\end{align*}
对客户端按 $q_j$ 加权求和，并使用 $\sum_j q_j=1$，即可得到 \eqref{eq:stale-bound}。
\end{proof}

\section{随机缓存联邦方法的收敛阶数}

令 $\widehat g^t_{\text{缓}}$ 表示实际用于原始变量更新的随机缓存方向。写成
\[
    \widehat g^t_{\text{缓}}
    =
    g(v^t)
    +
    \xi_t
    +
    b_t,
\]
其中
\[
    \E[\xi_t\mid\F_t]=0,
    \qquad
    \E[\norm{\xi_t}^2\mid\F_t]\le \Sigma_{K,B}^2,
    \qquad
    \norm{b_t}\le \Delta_t.
\]
由引理 \ref{lem:stale}，可取
\[
    \Delta_t
    =
    RY_\B\sum_{s=t-A}^{t-1}\sigma_s.
\]

\begin{lemma}[随机一步不等式]
\label{lem:stoch-one-step}
在假设 \ref{assump:bounded} 下，若原始变量更新使用 $\widehat g^t_{\text{缓}}$，
则对任意 $(\beta,v)\in\B\times\V$，存在绝对常数 $c_1>0$，使得
\begin{align}
    \E_t[
    \mathcal{L}(\beta^{t+1},v)
    -
    \mathcal{L}(\beta,v^{t+1})
    ]
    &\le
    \frac{1}{2\eta_t}
    \left(
        \norm{\beta-\beta^t}^2
        -
        \E_t\norm{\beta-\beta^{t+1}}^2
    \right)
    \nonumber\\
    &\quad+
    \frac{1}{2\sigma_t}
    \left(
        \norm{v-v^t}^2
        -
        \E_t\norm{v-v^{t+1}}^2
    \right)
    \nonumber\\
    &\quad+
    c_1\eta_t\Sigma_{K,B}^2
    +
    D_\beta\Delta_t,
    \label{eq:stoch-one-step}
\end{align}
其中 $\E_t[\cdot]=\E[\cdot\mid\F_t]$。
\end{lemma}

\begin{proof}
与确定性一步不等式相比，唯一差异是原始变量更新方向从 $g(v^t)$ 变成
$g(v^t)+\xi_t+b_t$，于是多出误差项
\[
    \langle \xi_t+b_t,\beta-\beta^{t+1}\rangle.
\]
对零均值噪声项，使用杨氏不等式：
\[
    \langle \xi_t,\beta-\beta^{t+1}\rangle
    \le
    \frac{\eta_t}{2}\norm{\xi_t}^2
    +
    \frac{1}{2\eta_t}\norm{\beta-\beta^{t+1}}^2.
\]
第二项可吸收到近端三点不等式中的平方距离项；取条件期望并用
$\E_t\norm{\xi_t}^2\le\Sigma_{K,B}^2$，得到 $c_1\eta_t\Sigma_{K,B}^2$。
对偏差项，利用 $\beta,\beta^{t+1}\in\B$：
\[
    \langle b_t,\beta-\beta^{t+1}\rangle
    \le
    \norm{b_t}\norm{\beta-\beta^{t+1}}
    \le
    D_\beta\Delta_t.
\]
代回确定性一步不等式即得 \eqref{eq:stoch-one-step}。
\end{proof}

\begin{theorem}[不用方差缩减时的显式 $\mathcal{O}(1/\sqrt{T})$ 阶]
\label{thm:main}
假设 \ref{assump:bounded}、\ref{assump:saddle} 和 \ref{assump:delay} 成立。令
\[
    \eta_t=\sigma_t=\frac{\gamma}{\sqrt{T}},
    \qquad
    t=0,\ldots,T-1,
\]
其中 $\gamma$ 足够小，使得原始--对偶混合梯度法的稳定条件成立。则随机缓存联邦盒对偶方法的遍历间隙满足
\begin{align}
    \E[
    \Gap(\bar{\beta}^T,\bar v^T)
    ]
    &\le
    \frac{1}{\sqrt{T}}
    \left[
        \frac{D_\beta^2+D_v^2}{2\gamma}
        +
        c_1\gamma
        \frac{c_0G^2}{K}
        \sum_{j=1}^m
        \frac{q_j^2}{\pi_j}
        \left(
            1+\frac{1}{B_j}
        \right)
    \right]
    \nonumber\\
    &\quad+
    \frac{D_\beta}{T}
    \sum_{t=0}^{T-1}
    RY_\B
    \sum_{s=t-A}^{t-1}\sigma_s ,
    \label{eq:main-rate}
\end{align}
其中 $D_v:=\sup_{v,v'\in\V}\norm{v-v'}$。进一步，由
$\sigma_s=\gamma/\sqrt{T}$ 可得
\begin{equation}
    \E[
    \Gap(\bar{\beta}^T,\bar v^T)
    ]
    \le
    \frac{1}{\sqrt{T}}
    \left[
        \frac{D_\beta^2+D_v^2}{2\gamma}
        +
        c_1\gamma
        \frac{c_0G^2}{K}
        \sum_{j=1}^m
        \frac{q_j^2}{\pi_j}
        \left(
            1+\frac{1}{B_j}
        \right)
        +
        D_\beta RY_\B A\gamma
    \right].
    \label{eq:sqrt-rate}
\end{equation}
因此，在 $K,B_j,A$ 固定且有界的情况下，不使用方差缩减也有
\[
    \E[
    \Gap(\bar{\beta}^T,\bar v^T)
    ]
    =
    \mathcal{O}(T^{-1/2}).
\]
\end{theorem}

\begin{proof}
记 $\eta=\gamma/\sqrt{T}$。在引理 \ref{lem:stoch-one-step} 中取
$(\beta,v)=(\beta^\star,v^\star)$，并对
$t=0,\ldots,T-1$ 求和。平方距离项望远镜相消，得到
\[
    \sum_{t=0}^{T-1}
    \E[
    \mathcal{L}(\beta^{t+1},v^\star)
    -
    \mathcal{L}(\beta^\star,v^{t+1})
    ]
    \le
    \frac{D_\beta^2+D_v^2}{2\eta}
    +
    c_1\sum_{t=0}^{T-1}\eta_t\Sigma_{K,B}^2
    +
    D_\beta\sum_{t=0}^{T-1}\Delta_t.
\]
利用 $\mathcal{L}$ 的凸凹性，将左侧转化为遍历平均点处的间隙，再除以 $T$。
代入引理 \ref{lem:variance} 的 $\Sigma_{K,B}^2$ 和引理 \ref{lem:stale} 的
$\Delta_t$，得到 \eqref{eq:main-rate}。最后代入
$\eta_t=\sigma_t=\gamma/\sqrt{T}$，并使用
\[
    \sum_{s=t-A}^{t-1}\sigma_s
    \le
    A\gamma/\sqrt{T},
\]
即可得到 \eqref{eq:sqrt-rate}。
\end{proof}

\section{固定步长情形}

\begin{corollary}[固定步长误差邻域]
若采用固定步长 $\eta_t=\sigma_t=\eta$，同样的推导给出
\begin{equation}
    \E[
    \Gap(\bar{\beta}^T,\bar v^T)
    ]
    \le
    \frac{D_\beta^2+D_v^2}{2\eta T}
    +
    c_1\eta\Sigma_{K,B}^2
    +
    D_\beta RY_\B A\eta.
    \label{eq:fixed-rate}
\end{equation}
因此固定步长有 $\mathcal{O}(1/T)$ 的暂态下降项，但最终会停在由随机方差和陈旧偏差决定的邻域中：
\[
    \mathcal{O}(\eta\Sigma_{K,B}^2)
    +
    \mathcal{O}(A\eta).
\]
\end{corollary}

\section{陈旧加权、稳健加权、自适应加权与方差缩减}

\paragraph{陈旧加权。}
如果服务器使用陈旧程度权重 $a_j^t$，并且
\[
    0<a_{\min}\le a_j^t\le a_{\max}<\infty,
    \qquad
    \sum_{j=1}^m q_j a_j^t=1,
\]
则上述证明保持成立，只是 $G$ 和陈旧偏差常数会乘上与 $a_{\max}/a_{\min}$ 有关的因子。

\paragraph{稳健加权。}
如果使用等客户端权重或自定义客户端权重，则目标从样本平均目标变为
\[
    F_w(\beta)
    =
    \sum_{j=1}^m
    w_j
    \frac{1}{n_j}
    \sum_{i\in\mathcal{D}_j}
    \rho_\tau(y_i-x_i^\top\beta)
    +
    r(\beta).
\]
同样的证明适用于 $F_w$ 对应的鞍点问题。这解释了为什么稳健版本可能牺牲样本平均目标值，却改善最差客户端损失。

\paragraph{自适应加权。}
若自适应权重 $w_j^t$ 满足有界性和有限总变化量
\[
    0<w_{\min}\le w_j^t\le w_{\max}<1,
    \qquad
    \sum_{t=0}^{T-1}\norm{w^{t+1}-w^t}_1<\infty,
\]
则可视为求解一个缓慢变化的加权鞍点问题序列。收敛界中会多出权重变化项
\[
    \frac{1}{T}
    \sum_{t=0}^{T-1}
    \norm{w^{t+1}-w^t}_1.
\]
若该项为 $\mathcal{O}(T^{-1/2})$ 或更小，则不改变主阶数。

\paragraph{方差缩减。}
方差缩减不是得到定理 \ref{thm:main} 的必要条件。若方差缩减后的方向满足
\[
    \E[
    \norm{\widehat g_{\rm vr}^t-g(v^t)}^2
    \mid\F_t]
    \le
    \Sigma_{\rm vr}^2
    \le
    \Sigma_{K,B}^2,
\]
则定理 \ref{thm:main} 仍成立，只需将 $\Sigma_{K,B}^2$ 替换为更小的
$\Sigma_{\rm vr}^2$。因此，方差缩减改善的是收敛界中的方差常数，而不是在一般非光滑随机联邦设定下把阶数从 $\mathcal{O}(T^{-1/2})$ 直接提升为 $\mathcal{O}(T^{-1})$。

\section{结论}

在明确的有界性、客户端抽样、小批量抽样和有界延迟假设下，我们得到：
\begin{itemize}[leftmargin=2em]
    \item 全量确定性盒对偶原始--对偶混合梯度法：
    \[
        \Gap(\bar{\beta}^T,\bar v^T)=\mathcal{O}(1/T).
    \]
    \item 随机联邦缓存盒对偶方法，不使用方差缩减：
    \[
        \E[\Gap(\bar{\beta}^T,\bar v^T)]
        =
        \mathcal{O}(1/\sqrt{T}).
    \]
    \item 显式常数依赖客户端参与数 $K$、小批量大小 $B_j$、客户端不均衡因子
    $\sum_j q_j^2/\pi_j$、最大缓存延迟 $A$、协变量上界 $R$ 和残差上界
    $Y_\B$。
    \item 固定步长时，方法收敛到由随机方差和陈旧偏差决定的误差邻域。
    \item 方差缩减可以降低方差常数，但不是证明 $\mathcal{O}(1/\sqrt{T})$ 阶的必要条件。
\end{itemize}

\end{document}
